\documentclass{PASJ2026}
\usepackage{amsmath,amssymb}
%\usepackage[pdfencoding=auto]{hyperref}
\usepackage{pifont}
% \usepackage{flushend} % balance the last page
% (in order to balance the last page you can use one of \balance or \usepackage{flushend})
%
% This source file can be processed by LuaTeX or upTeX.
% for LuaTeX
%  latexmk -lualatex TeXtemplate2026.tex
% for upTeX
%  latexmk -pdfdvi TeXtemplate2026.tex
%  (latexmkrc is needed in the same directory.)
%
% versions for LuaTeX and upTeX were merged. K.F
\ifLuaTeX  % LuaTeX
\usepackage{luatexja-fontspec}
\setmainfont{TimesNewRoman}
\setsansfont{Arial}
\defaultjfontfeatures{Scale=1.0}
% we should use the following line in the future, K.F
% on MacOS at least, Hiragino font seems to be better than Ipaex
\setmainjfont[BoldFont=IPAexGothic]{IPAexMincho}
\setsansjfont{IPAexGothic}
% other possibilities
%\setmainjfont[BoldFont=HiraMinProN-W6]{HiraMinProN-W3}
%\setsansjfont[BoldFont=HiraginoSans-W8]{HiraginoSans-W4}
%\setmainjfont[BoldFont=IPAexGothic]{KozMinPr6N-Light}
%\setmainjfont[BoldFont=IPAexGothic]{MS-PMincho}
%\setsansjfont{MS-PGothic}
% \else  % upTeX
% for upLaTeX default fonts are employed
\fi  % LuaTeX/upTeX
%%
%%
%%
%%
%%
%%
%%
%
% 図キャプション表タイトルの上下の行間隔を調整
\abovecaptionskip=3pt
\belowcaptionskip=3pt
% 所属機関記号 Affiliation
\newcommand{\jaffi}[1]{$^{\mathrm{#1)}}$}
\newcommand{\affi}[1]{$^{\,\mathrm{#1)}}$}
\newcommand{\affithanks}[2]{\,\thanks{#1}\!$^{,\mathrm{#2)}}$}
%
\begin{document}
%
\jtitle{日本加速器学会年会・論文投稿要領}
%
\title{Preparation of Papers for Accelerators Conferences}
%
\jauthor
{
丸井花子\jaffi{*, A},
真直太郎\jaffi{A},
加速粒子\jaffi{B}
}
%
\author
{
Hanako Marui\affithanks{pasj-desk@conf.bunken.co.jp}{A},
Taro Masugu\affi{A},
Ryushi Kasoku\affi{B}\\
\affi{A} One of Accelerator Facilities\\
\affi{B} Office of Particle Accelerator Society of Japan\\
}
% 所属記号を使用しない場合は、\affithanks{Email}{記号} の代わりに \thanks{Email} を使用して下さい
%
\Abstract
{
This template has been prepared for writing papers in Japanese. Please use the JACoW template to write your English paper. If you write texts directly into this .tex file, you can easily complete your paper with \TeX. When converting your completed paper to PDF, be sure to embed all fonts. The number of pages in the paper should be 2-5, excluding references. Please note that the paper style for PASJ2026 has been changed as follows. 1) The title and author's name should be written in Japanese and English. 2) Figures and tables and their captions should be written in English. When creating English abstracts and figure captions using translation software, please refer to the notes in the text below. (Translated by Google)
}
%
\maketitle
%
\section{はじめに}
論文執筆にあたっては、必ず最新の年会プロシーディングスのテンプレートを使用してください。%このテンプレートを使うと、メニューからスタイルを開いてスタイルリストから選択するだけで、原稿のカーソル位置にスタイルが設定できます。
このテンプレートを使わないと適切なスタイルが設定できずに、編集作業に非常に手間がかかってしまう場合があります。

PASJ2023以降のテンプレートでは論文のフォーマットの修正のほか、いくつかの変更点や注意点の追記があります。特に以下の要点に十分ご注意ください。
%
\begin{Itemize}
\item{タイトルおよび著者名は日本語と英語を併記、日本語英語とも著者名の区切りは半角カンマ（その後ろに半角スペース1個）、所属参照記号は全て半角Times系フォントで上付き。} 
%
\item{日本語の文章は句読点（「。」と「、」）を使って記述。}
% 公用文作成の考え方
% https://www.bunka.go.jp/seisaku/kokugo_nihongo/kokugo_shisaku/94336802.html

%
\item{数値と単位記号の間は詰めず、通常の半角スペースではなくnon-breaking spaceを挿入。\LaTeX ではできれば 1/4スペース、$\backslash${,} または \textasciitilde$\backslash${!} が好ましい。}
% この目的の siunitx package は、残念ながら uplatex と graphicsx とを組み合わせると問題を起こすので、lualatex でしか使えない。
%
\item{図表とその説明文は英語で記述。図や写真はFigure（短縮はFig.）、表はTableと記して参照。}
%
\end{Itemize}

なお、翻訳ソフトを使用して要旨や図表の英文を作成する場合は、\ref{honyaku}を参考にしてください。

その他、以下についてもご注意お願い致します。
%
\begin{Itemize}
% \item{PDFに変換する際は、必ず全てのフォントを埋め込むこと。PDFはバージョン1.4で作成。} 
%
\item{原稿枚数は、2$\sim$5ページです。書式指定に従い執筆してください。なお、原稿末尾の参考文献リストのみ、より正確な記述のために6ページ目以降にまたがっても結構です。}
%
\end{Itemize}

論文は、英語か日本語で記述していただきます。英語の原稿の場合は、JACoWの書式に従ってください。この際、タイトル・著者名・所属に日本語の記述は不要です。JACoW形式のテンプレートはJACoWのホームページ~\cite{Reference_JACoW_Word,Reference_JACoW_TeX}からダウンロードできます。

以下の説明はすべて日本語で原稿を記述する場合のものです。日本語で記述する場合も、JACoWの書式を日本語化したもの~\cite{Reference_1,Reference_1_URL,T_Tanaka}に準じていますが、Webでの公開に伴い英語による検索・参照に対応できるよう、論文題名・著者名・所属・要旨・図表・図表番号及び説明文は必ず英語表記を入れてください。

以上の理由により、日本語原稿の書式はJACoWとは若干異なります。以下の説明文書を熟読の上、正しい書式の原稿を完成させるようご協力ください。
その他、ご不明な点等は加速器学会事務局（pasj-desk@conf.bunken.co.jp）にお尋ねください。
%
\section{本テンプレートを使った原稿の書き方}
本{\TeX}用テンプレートは、羽島良一氏（日本原子力研究開発機構）がJACoWの{\TeX}用テンプレートを元に
2004年に作成したリニアック技術研究会論文用{\TeX}ファイル~\cite{Reference_2}及びクラスファイルを、
加速器学会の書式に合うように修正したものです。

本テンプレートは、本文が記載された{\TeX}ファイルTeXtemplate2026.tex、および書式の細部が記述されたクラスファイルPASJ2026.clsそして画像ファイルJACpic\_mc.pdf, chart.png, chart2.jpgから成っています。
論文執筆にあたっては、TeXtemplate2026.texファイルに直接書き込むか、別に用意したテキストを挿入するのが最も簡単です。
本テンプレートの文字エンコーディングはUTF-8、改行コードはCR+LFです。これらを変更されたい方はmi~\cite{mi}等のテキストエディタを使用して変更して下さい。

原稿は、2$\sim$5ページの範囲で書式指定に従って執筆してください。ただし、原稿末尾の参考文献リストについては6ページ目以降にまたがっても結構です。

書式の詳細や書き方の注意は、第\ref{kakikata}章に記載されています。
また、JACoWのwebページ~\cite{Reference_3}には、論文執筆に役に立つ情報が多く記載されていますので、是非参考にしてください。
%
\section{原稿の書式と書き方}
\label{kakikata}
%
\subsection{用紙の設定}
用紙の大きさはA4、本文は二段組（2カラム）とします。
余白等の詳細はTable~\ref{Tab1}およびFig.~\ref{Fig1}の通りです。
%
\begin{table}[hbt]
\centering
\caption{Document Layout}
\begin{tabsmall}
\begin{tabular}{ll}
\toprule
Paper size & A4 (210~mm $\times$ 297~mm) \\
Margin & Top : 37~mm, Bottom : 19~mm \\
 & Left : 20~mm, Right : 20~mm \\
Column & 2 columns\\
 & Width : 82.5~mm, Spacing : 5.0~mm \\
\bottomrule
\end{tabular}
\end{tabsmall}
\label{Tab1}
\end{table}
%
\begin{figure}[hbt]
\centering
\includegraphics[width=55mm]{JACpic_mc.pdf}
\caption{Document layout.}
\label{Fig1}
\end{figure}

\subsection{スタイルと使用フォント}
フォントは、指定のない限り図も含めて、できるかぎり明朝系、ゴシック系（日本語）、Times系（英数字、記号、特殊文字）に限定してお使いください。
{Lua\TeX}では$\backslash$usepackage\{luatexja-fontspec\}でフォントを指定出来ます。本テンプレートでは日本語メインフォントをIPAex明朝（但し太字はIPAexゴシック）
日本語サンセリフフォントをIPAexゴシックとしていますが、フォントがPC内に存在しない場合は他のフォントを使用して下さい。

なお、日本語{\TeX}のフォントファミリは明朝系（セリフ/メイン）及びゴシック系（サンセリフ）の2つで、フォントシリーズは標準（Medium）のみ、太字（Boldface）はありません。
英数字のフォントファミリはTimes系を使ってください。$\alpha\beta\gamma$等のギリシャ文字には数式モード\$$\backslash$alpha$\backslash$beta$\backslash$gamma\$を使用して下さい。

本テンプレートの書式はMicrosoft Word用に定義されたスタイルに準じていますが、完全に同じでは有りません。スタイル毎のフォントのサイズや行間については、Table~\ref{Tab2}にまとめられています。

PDFに出力した時の1行の高さがWordとほぼ同じになるように、クラスファイル中でフォントサイズ Large, large, normalsize, tabsmallなどを再定義しています。
%
\begin{table*}[hbt]
\centering
\caption{Definition of Text Styles}
\begin{tabsmall}
\begin{tabular}{lccccc}
\toprule
 Style & Fontfamily &  Fontseries & Fontsize & Size name & Baselineskip\\
\midrule
Japanese title & Sans(Gothic) & Medium & 14~pt & Large & 19.6~pt\\
TITLE & Main(Times) & {\textbf{Boldface}} & 12~pt & large & 16.7~pt\\
Japanese author & Main(Mincho) & Medium & 10~pt & authorf & 16~pt\\
Author & Main(Times) & Medium & 10~pt & authorf & 16~pt\\
Section title & Sans(Gothic) & Medium & 12~pt & large & 16.7~pt\\
Subsection title & Main(Mincho) & Medium & 10~pt & normalsize & 14~pt\\
Text & Main(Mincho) & Medium & 10~pt & normalsize & 14~pt\\
Table title & Main(Times) & Medium & 10~pt & normalsize & 12~pt\\ 
Figure caption & Main(Times) & Medium & 10~pt & normalsize & 12~pt\\ 
Table & Main(Times) & Medium & 9~pt & tabsmall & 16~pt\\
References & Main(Mincho) & Medium & 9~pt & small & 11~pt\\
References(URL,doi) & Mono(LMMono) & Medium & 9~pt & small & 11~pt\\
\bottomrule
\end{tabular}
\end{tabsmall}
\label{Tab2}
\end{table*}
%
\subsection{論文題名・著者・要旨}
論文題名および著者名は、日本語と英語の併記となりますので十分ご注意ください。
ただし、著者所属および要旨は従来通り、英語のみです。

書式は以下の通りです。

\begin{Itemize}
\item 日本語題名（英数字含む）はゴシック14~pt、すぐ次の行に英語題名を全て大文字でTimes系太字12~pt。
物理単位の小文字や慣習上の小文字は小文字のまま。
%
\item 日本語著者名、英語著者名、英語所属機関名の順に、それぞれを改行で区切り、頭文字のみ大文字。
フォントサイズは10~pt。
著者名の区切りは日本語、英語とも、半角カンマ＋半角スペース1個。
%
\item 所属機関の所在地は不要。
%
\item 所属機関が複数の場合には、A, B, $\cdots$のように半角Times系フォントで上付参照記号を振る。
%
\item 代表者の電子メールアドレスを脚注に記載する場合、例のように上付参照記号 $*$, $\dagger$, $\ddagger$ 等を付ける。
%
\item 要旨はイタリック12~ptで{\textit Abstract}と見出しをつける。Abstract文は改行なしの一段落のみで、英語10~ptで記述。
%
\end{Itemize}
%
\subsection{本文}
% 本文の始まる位置（上から測った高さ）を指定することもできます。
% このTeXtemplate2026.texファイルプリアンブル中の$\backslash$setlength\{$\backslash$titleblockheight\}を指定してください。

本文は二段組で記述します。フォントは、10~ptの明朝系を使用し、必要に応じて、{\textbf{ゴシック系}}を適宜使用してください。
ただし英数字はTimes系を使用してください。
数値と単位記号の間にはnon-breaking space（半角のチルダ記号 ``\~'' ）を挿入してください。
%
\subsection{章・節・項の見出し}
各章の見出し、節の小見出しには必ず番号を付けてください。見出しのフォントはゴシック系12~ptを、小見出しは明朝系10~ptをお使いください。

なお、本テンプレートでは、見出し、小見出しに自動的に番号を振るようになっています。
%
\subsection{箇条書き}
箇条書きの部分は、$\backslash$begin\{Itemize\}および$\backslash$end\{Itemize\}ではさんでください。

Itemizeをitemizeとすると、各項目の間隔が広くなります。
%
\subsection{図表の挿入の仕方}
図および表とそれらの説明文（キャプションおよびタイトル）は英語で記述してください。

作成にあたっての注意点は以下の通りです。
%
\begin{Itemize}
%
\item{図および表の説明文の始まりを、図は「Figure~\ref{Fig1}:」、表は「Table~\ref{Tab1}:」などとし（数字の前とコロンの後に半角スペースを入れる）、「Fig.~\ref{Fig1}:」等略記は不可。}
%
\item{表のタイトルの記述にあたってはキャピタライゼーション・ルールに関する解説（\cite{Capitalization_Rules}のTitlesの項目）を参考に、\\
%1) タイトルの最初と最後の単語の頭文字を大\\
%　文字\\
1) タイトルの最初と最後の単語の頭文字を大文字\\
2) 名詞、代名詞、動詞、形容詞、従属接続詞、副詞\\
　の頭文字は大文字\\
3) 冠詞、等位接続詞、前置詞、不定詞のtoは小文\\
　字\\
を基本とする。}
%
\item{Table~\ref{Tab1}の例では、表本体の部分を$\backslash$begin\{tabsmall\}および$\backslash$end\{tabsmall\}ではさみ、
フォントサイズを小さめ（9~pt）、1行を広め（16~pt）に指定している。}
%
\item{図のキャプションの文末にはピリオドが必要。}
%
\item{Figure~\ref{Fig1}の例では、PDFを幅55~mmに貼り込んでいる。
この例では、$\backslash$begin\{figure\}の位置指定オプションにhbtを付け、「その位置」への挿入を優先するよう指定している。}
%
\item{Figure~\ref{Fig2}の例では、subfloatを定義して図を二枚並べている。図の位置はtbpとして上か下、あるいは次のページを指定している。}
%
\item{Figure~\ref{Fig2}の例では、PNGとJPEGファイルを各幅80~mmに貼り込んでいる。}
%
\item{横二段に渡る図を入れる場合は、$\backslash$begin\{figure$*$\}および
$\backslash$end\{figure$*$\}のようにfigureの次に$*$を付けて指定する。表についても
同様に$\backslash$begin\{table$*$\}および$\backslash$end\{table$*$\}とする。
ただし、1ページ目に横二段の図表を挿入することは出来ないので注意。}
%
\item{図および表の説明文の配置は、それが1行に収まる場合は段の中央寄せ、2行以上は段幅（82.5~mm）一杯の両端揃え。
なお、本TeXtemplate2026.texファイルプリアンブル中の$\backslash$abovecaptionskip及び$\backslash$belowcaptionskipで、
図キャプション表タイトルの上下それぞれに3~ptの空間を設定している。}
%
\end{Itemize}
%
\begin{figure*}[tbp]
\centering
\subfloat[PNG image (430 $\times$ 200)]
{
  \includegraphics[width=80mm]{chart.png}
  \label{fig:two_a}
}
\subfloat[JPEG image (430 $\times$ 200)]
{
  \includegraphics[width=80mm]{chart2.jpg}
  \label{fig:two_b}
}
\caption
{
These examples are the circular charts showing the shares of PC's OSs used to write the papers for IPAC10.
(Publishing approved by JACoW)
}
\label{Fig2}
\end{figure*}
%
\subsection{図作成における注意}
%
\begin{Itemize}
%
\item 図のフォーマットはビットマップが一番確実で、以下のように選ぶと美しさとサイズのバランスが最適。
%
\begin{description}
\item[] 写真：JPEG
\item[] CAD図・グラフ等：PNG（TIFFも可）
\end{description}
%
\item PNGはGIFを拡張した優れたビットマップフォーマットで、線画等を劣化無くコンパクトに圧縮できるため、強く推奨する。
Figure~\ref{Fig2}は、円グラフをPNGとJPEGに変換した例で、いずれも約12~kBのサイズであるが、JPEGでは画質が大きく低下していることがわかる。
%
\item 写真やグラフは、そのファイルサイズが大き過ぎると表示に時間がかかってしまうので、なるべく小さなファイルに変換し挿入すること。特に EPS や WMF / EMF は、作り方によってはとても大きくなることがあるので注意が必要。
%
\item 論文を原寸で印刷する場合、PNGやJPEGの図の分解能は、特に文字が小さくない限り最低300~dpi有ればよく、
この場合一段分の幅は約1000ドットに相当する。
%
\item 多数の点をプロットしたグラフをメタファイルとして貼り付けると、論文のデータサイズが大きくなり過ぎるため、データを間引くか上記のようにPNG 等のビットマップに変換する等して、データサイズを小さくすること。
%
\item 図表はカラーでも良いが、モノクロプリンタで印刷する場合でも読者に理解できるよう、配色に注意する。
%
\item 図中も原則として所定のフォントを使用する。
%
\end{Itemize}
%
\subsection{数式の挿入}
数式は、{\TeX}のコマンドを使って作成し、必要なら数式番号を(\ref{equation_1})のように振ってください。
\begin{equation}
E = \frac{m_0 c^2}{\sqrt{1 - \beta^2}}
\label{equation_1}
\end{equation}
%
\subsection{図と表の参照}
本文中で図や表を参照するときには、Tableの場合は全てTable~\ref{Tab1}
（数字の前に半角スペースを入れる）などとし、
一方図の場合には、文頭にあるときはFigure~\ref{Fig1}
（数字の前に半角スペースを入れる）のように略さず記し、
文中ではFig.~\ref{Fig1}（数字の前に半角スペースを入れる）のように略記してください。
%
\subsection{数式の参照}
本文中で数式を参照するときには、文頭ではEquation (1)（左カッコの前に半角スペースを入れる）、
文中ではEq. (1)（左カッコの前に半角スペースを入れる）のように略記してください。
%
\subsection{文献・URLの参照}
記事に関連する説明等を補足するために文献及びURLを参照する事をおすすめします。本文中で参照する文献及びURLには、文中で初めて引用される順に通し番号を振り、 \textbackslash{cite}\{\}を使用してください。

同じ位置で複数の文献を参照するときには、\cite{Reference_JACoW_Word, Reference_JACoW_TeX}、\cite{Reference_JACoW_Word, Reference_JACoW_TeX, Reference_1}、\cite{Reference_JACoW_Word, Reference_JACoW_TeX, Reference_1, T_Tanaka}、\cite{Reference_JACoW_Word, Reference_JACoW_TeX, Reference_1, T_Tanaka, Reference_2}、\cite{Reference_JACoW_Word, Reference_JACoW_TeX, Reference_1, T_Tanaka, Reference_2, mi}のように記述することを推奨します。
%
\subsection{参考文献リストの書式}
\subsubsection{参考文献の書式}
\label{reference_format}
インデントを揃えるために、参考文献の数が9編以下ならば$\backslash$begin\{thebibliography\}\{9\}を、10編以上ならば$\backslash$begin\{thebibliography\}\{99\}を選択してください。

参考文献は、9~ptで、日本語は明朝系、英文の場合はTimes系フォントで記載してください。

文献のみの記述の場合~\cite{Reference_1}、著者名（複数の場合カンマ+半角スペースで区切り、3名を超える場合は最初の著者以外はまとめてカンマを入れず斜体で \textit{et al.} として良い）、``著書・論文タイトル''、紙誌名、出版年、ページ範囲、の順を基本とし、最後はピリオド「.」で終わります。なお、\textit{et al.} 以外の文字は原則として斜体にはせず、太文字の使用は不要です。

\subsubsection{DOIの記述}
参考文献が電子出版されていて、DOIで参照できる場合にはDOIを追記することを推奨します。

DOIは、参考文献の最後のピリオドの後に半角スペースまたは必要であればインデント付き改行を挿入し、その後\textbackslash{url}\{\}を使用してdoi:またはDOI:からスペースなしで1行に納まるよう記述してください。DOIの最後にはピリオド等を入れないでください。DOIの記述に関してはJACoWのテンプレート~\cite{Reference_JACoW_TeX}が参考になります。

\subsubsection{URLの記述}
電子出版されている参考文献のURLを追記する場合、\ref{reference_format}の書式に従って参考文献を記述した最後のピリオドの後に半角スペースを挿入し、続けてURLを記述することとし、最後にはピリオド等を入れないでください。

参考文献がWebコンテンツである場合~\cite{Reference_2, item:2-1, item:7, item:8}、なるべくURLの前にWebページのタイトルを記入して下さい。
その場合、Webページタイトルの後にカンマと半角スペースを挿入して区切り、その後\textbackslash{url}\{\}を使用してURLを続けて記述することとし、最後には\textbf{ピリオド等を入れない}でください。URLのみ記述の場合~\cite{Reference_1_URL}も同様。

URLが長過ぎて1行に収まらない時は、途中にスペースを入れるなど2行にわたるようにしてください。

\subsubsection{参考文献の表記の仕方}
様々な参考文献の表記の仕方について厳密に指定してはいませんが、読者が容易に参照可能となるよう十分な情報の記述を心がけてください。この際、JACoWのテンプレート~\cite{Reference_JACoW_TeX}におけるReferenceの表記の例が参考になります。

なお、日本加速器学会年会のプロシーディングスを参照する際には年会名称を~\cite{item:9,item:10}の例のように記述するほか、~\cite{item:11,item:12}の略記のみでも結構です。ただし、PASJ\#\#\#\#と正式に西暦の略記が用いられているのは2014年（第11回年会）以降の年会に限られていますので、ご承知ください。

また、当該年会での発表を参照する場合、未出版であることからプロシーディングスを参照するのではなく~\cite{item:13,item:14}のように年会名称と講演番号を記述し、単に ``this meeting'' のみのような出展が具体的ではない記述は避けてください。
%
\subsection{最終ページの原稿の左右2段均等配分}
最終ページの原稿の左段（またはそれより前）で\textbackslash{balance}を一回実行してください。これにより最終ページの原稿の体裁を、原則として左右2段に均等配分します。このとき、原稿の行数や挿入されている図表のサイズやキャプションの長さ等が原因で、左右2段均等配分からずれる場合が生じるのは構いません。
%
\subsection{脚注の挿入}
脚注を挿入する場合は、\textbackslash{footnote}\{\}を使用してください。
本文中に挿入すれば\footnote{脚注は\textbackslash{footnote}\{\}を使用して記入。}、自動的に番号が振られ、段ごとにまとめてページ最下段に配置されます。
%
\subsection{自動翻訳利用の際の注意}
\label{honyaku}
要旨や図の日本語説明文をウェブサイトの自動翻訳を利用して英文に変換する場合、いくつか注意をすることでかなり英文を理解しやすくできます。その注意を以下に列挙します。
%
\begin{enumerate}
%
\item{技術論文の場合、一般的なサイトの中では「Google翻訳」や「DeepL」翻訳などは比較的良好な翻訳結果が得られるようである。また、最近は対話型生成AIサービスも実用上十分な翻訳結果を提供することが多い。}
%
\item{自動翻訳システムが誤解なく理解しやすい文を入力すること。そのような文を書くに当たっての注意は以下の通り。
なお、一般にそのような文は、我々から見ればくどくて不自然である。}
\begin{Itemize}
%
\item{一つの文は出来るだけ短くする。}
%
\item{主語を省略せず明確にする。}
%
\item{省略した表現や名称を使わない。}
%
\item{専門用語等は初めから英語にしておく。}
%
\item{名詞を他の名詞の形容詞として使う場合、語順や名詞間にいれる「の」等の言葉により翻訳結果はかなり異なる場合がある。}
%
\item{くどくても、未来や過去など時制が明解な表現とする。}
%
\end{Itemize}
%
\item{いかに上記の注意を払っても、不自然あるいは不正確な英訳をなくすことは出来ない。
むしろ、著者が上記の観点からすぐに不適切と分かる日本語をまず修正してから自動翻訳し、
その後英訳を修正する方が早い。}
%
\item{特に日本語にはない冠詞や複数形は自動翻訳では不正確なので、著者が修正する必要がある。}
%
\end{enumerate}

このテンプレートのAbstractは、以下の文章をGoogleの翻訳ページで英語に変換したものです。
日本語の表現を若干工夫して、なるべく理解できる英語になるようにしてありますが、翻訳された英語は一切修正していません。

【Abstract翻訳元の日本語要旨】
\begin{description}
\item[]
{
　本テンプレートは論文を日本語で記述するために用意されました。英語論文執筆には JACoW のテンプレートをお使いください。あなたがこの.texファイルに直接textsを書き込めば、{\TeX}で論文を容易に完成させることができます。完成した論文をPDFに変換する際は、必ず全てのフォントを埋め込んでください。論文のページ数は、参考文献を除いて2-5ページでなければならない。PASJ2026の論文スタイルが次のように変更されましたので、ご注意ください。1) タイトルおよび著者名は日本語と英語で書くこと。2) Figures and tablesとそれらのcaptionsは英語で記述すること。翻訳ソフトを使用して英語の要旨やfigure captionsを作成する場合は、以下の本文中の注意事項を参考にしてください。%本テンプレートは論文を日本語で記述するために用意されました。英語論文執筆には JACoW のテンプレートをお使いください。
}
\end{description}
%
\section{原稿の仕上げと投稿}
%
\subsection{チェック項目}
%
\begin{Itemize}
%
\item 代表者メールアドレスが正しく脚注にあるか。
%
\item ページ番号やヘッダー、フッターを入れない。
%
\item 図のキャプション、テーブルのタイトルの記述。

・1行のみでは中央寄せとし、複数行では段幅一杯に合わせる

・テーブルタイトルは頭文字を大文字とし、文末のピリオドなし（3行以上の長文は通常文章スタイルも可）

・図のキャプションは通常の文章スタイルでピリオド付き
%
\item 章や図表を参照する番号は正しいか。
%
\item 参考文献番号は正しいか。
%
\item シンボルがアルファベットに化けていないか。
%
\item 原稿のデータサイズが大き過ぎないか。（目安として10~MB以下）
%
\end{Itemize}
%
\subsection{PDFファイルの作成}
PDFファイルの作成は{\TeX}アプリケーションから直接行って下さい。% PDFはバージョン1.4で作成してください。
Macintosh及びWindowsのOSを使用されている方は、TeX Live\cite{MacTeX,WindowsTeX}の使用をおすすめします。
また、Cloud LaTeX~\cite{Cloudlatex}やOverleaf~\cite{Overleaf}といったウェブサービスも容易に利用ができ、無償での利用も可能です。
% テンプレートは2種類（TeXtemplate2026\_LuaTeX.tex及びTeXtemplate2026\_upTeX.tex）あり、クラスファイルも2種類（PASJ2026\_LuaTeX.cls及びPASJ2026\_upTeX.cls）あります。
%
\subsubsection{MacOS且つLuaLaTeXの場合}
TeXShop $\rightarrow$環境設定$\rightarrow$書類$\rightarrow$設定プロファイル$\rightarrow$ upTeX(ptex2pdf)を選び、
タイプセット$\rightarrow$ pdfTeXを選んで下さい。
\subsubsection{MacOS且つupLaTeXの場合}
TeXShop $\rightarrow$環境設定$\rightarrow$書類$\rightarrow$設定プロファイル$\rightarrow$ upTeX(ptex2pdf)を選び、
タイプセット$\rightarrow$ TeX + DVIを選んで下さい。
\subsubsection{WindowsOS且つLuaLaTeXの場合}
TeXworks $\rightarrow$タイプセット$\rightarrow$ LuaLaTeXを選んで下さい。
\subsubsection{WindowsOS且つupLaTeXの場合}
TeXworks $\rightarrow$タイプセット$\rightarrow$ upLaTeXを選んで下さい。
%
\subsubsection{OverleafまたはCloud LaTeX且つLuaLaTeXの場合}
MenuまたはProject Settingsで、CompilerまたはLaTeX Engineとして、LuaLaTeXを選んで下さい。\\
%
\subsubsection{Cloud LaTeX且つupLaTeXの場合}
Project Settingsで、LaTeX Engineとして、upLaTeXを選んで下さい。
\subsubsection{Overleaf且つupLaTeXの場合}
% Menuで、Compilerとして、LaTeXを選び、参考情報~\cite{Overleaf-uptex}に従いlatexmkrcという設定ファイルを作成して下さい。latexmkrcは本テンプレートにも含まれています。
Menuで、Compilerとして、LaTeXを選び、下に記載するlatexmkrcという設定ファイルをソースファイルと同じ場所に置いて下さい。

{up\LaTeX} を Engine に使用する場合には、latexmkrcという設定ファイルを用いて latexmk に LaTeX Engine を教えます。テンプレートには参考情報~\cite{Overleaf-uptex}を元にした次のような latexmkrc 設定ファイルが含まれていますので、ソースファイルと同じフォルダに作成して下さい。
\begin{description}
\item[latexmkrc の内容]{
\ \\
\$latex = `uplatex';\\
\$bibtex = `upbibtex';\\
\$dvipdf = `dvipdfmx \%O -o \%D \%S';\\
}
\end{description}

PDF変換の際は、必ず全てのフォントを埋め込むように変換してください。こうすることで、どのPCの上でも著者の意図した通りの原稿を再現することが出来ます。PDFファイルにフォントが埋め込まれているかどうかは、Adobe Acrobat Reader DCなどで文書のプロパティを見れば確認できます。

{Lua\LaTeX}と{up\LaTeX}の双方とも通常の論文では処理結果にほとんど差はありませんが、{Lua\LaTeX}は直接PDFファイルが生成できるので、PDFのきめ細かい調整が可能になります。{up\LaTeX}は日本語向けに開発されたので、日本語向けのきめ細かい調整が可能で、また処理速度が数倍高速になっています。

% 
\balance % 最終ページの左右段の均等配置

%
\subsection{投稿}
原稿の{\TeX}ファイル、図のファイル及び完成したPDFファイルを、定められた期日までに学会Webページのプロシーディングスファイルアップロード用ページからアップロードしてください。
%
\acknow{謝辞}
謝辞の見出しには章番号を振らないため、謝辞を入れる場合は、$\backslash$acknow\{謝辞\}と書き込み、謝辞の文章を続けて下さい。

%
%\begin{thebibliography}{9}
\begin{thebibliography}{99}

\bibitem{Reference_JACoW_Word}
MS Word templates of JACoW for Proceedings,
\url{http://www.jacow.org/Authors/MSWord}

\bibitem{Reference_JACoW_TeX}
LaTeX templates of JACoW for Proceedings,
\url{http://www.jacow.org/Authors/LaTeX}

\bibitem{Reference_1}
H.~Kobayashi~\textit{et al}., ``第26回リニアック技術研究会論文投稿要領'', Proceedings of the 26th Linear Accelerator Meeting in Japan, Tsukuba, Japan, Aug. 2001, pp. 0-0.

\bibitem{Reference_1_URL}
\url{http://conference.kek.jp/LAM26/lam26template.pdf}

\bibitem{T_Tanaka}
T.~Tanaka~\textit{et al}., ``日本加速器学会年会 論文投稿要領'', Proceedings of the 14th Annual Meeting of Particle Accelerator Society of Japan, Sapporo, Japan, Aug. 1-3, 2017, pp. 0-0. \url{http://www.pasj.jp/templates.html}

\bibitem{Reference_2}
R.~Hajima, ``Emittance Compensation in an ERL Merger'', Proceedings of the 29th Linear Accelerator Meeting in Japan, Funabashi, Aug. 2004, pp. 432-434. \url{http://www.pasj.jp/web_publish/pasj1_lam29/proceedings.html}

\bibitem{mi}
\url{http://www.mimikaki.net/}

\bibitem{Reference_3}
Electronic Processing of Submissions, \url{http://www.jacow.org/index.php?n=Authors.HomePage}

\bibitem{Capitalization_Rules}
\url{https://en.wikipedia.org/wiki/Capitalization}

\bibitem{item:2-1}
A.~Alpha and B.~T.~Beta, ``Novel techniques for future TeV electron accelerators'',
in \textit{Proc.\ IPAC'18}, 	Vancouver, BC, Canada, Apr.-May 2018,
pp.\ 567-569. %\newline
\url{doi:10.18429/JACoW-IPAC2018-PAPERID}
\bibitem{item:7}
P.~Venus \emph{et al.},
``New techniques in laser wakefield accelerators'',
\emph{Phys.\ Rev.\ ST Accel.\ Beams}, vol.\ 18,
p.\ 120198, 2015. %\newline
\url{doi:10.1103/PhysRevAccelBeams.18.120198} 
\bibitem{item:8}
T.~Earth \emph{et al.},
``Low dose irradiation impact on modern silicon detectors'',
\emph{Nucl.\ Instr.\ Meth.}, vol.\ 692, pp.\ 256--280, 2014.
\url{doi:10.1016/j.nima.2014.11.022}
\bibitem{item:9}
A. Sato \emph{et al.}, ``Virtual sample of reference description I'', Proc. 16th Annual Meeting of Particle Accelerator Society of Japan (PASJ2019), Kyoto, Japan, Jul.-Aug. 2019, pp. 10-13.
\bibitem{item:10}
I. Suzuki \emph{et al.}, ``Virtual sample of reference description II'', Proc. 19th Annual Meeting of Particle Accelerator Society of Japan (PASJ2022), Kitakyusyu (Online meeting), Japan, Oct. 2022, pp. 14-16.
\bibitem{item:11}
A. Sato \emph{et al.}, ``Virtual sample of reference description I'', Proc. PASJ2019, Kyoto, Japan, Jul.-Aug. 2019, pp. 10-13.
\bibitem{item:12}
I. Suzuki \emph{et al.}, ``Virtual sample of reference description II'', Proc. PASJ2022, Kitakyusyu (Online meeting), Japan, Oct. 2022, pp. 14-16.
\bibitem{item:13}
E. Sasaki \emph{et al.}, ``Virtual sample of reference description III'', The 15th Annual Meeting of Particle Accelerator Society of Japan (PASJ2018), Nagaoka, Japan, Aug. 2018, THP011, this meeting.
\bibitem{item:14}
O. Takahashi \emph{et al.}, ``Virtual sample of reference description IV'', PASJ2017, Sapporo, Japan, Aug. 2017, TUP123, this meeting.

\bibitem{MacTeX}
\url{https://tug.org/mactex/}
\bibitem{WindowsTeX}
\url{https://www.tug.org/texlive/acquire-netinstall.html}
\bibitem{Cloudlatex}
\url{https://cloudlatex.io/}
\bibitem{Overleaf}
\url{https://overleaf.com/}
\bibitem{Overleaf-uptex}
\url{http://doratex.hatenablog.jp/entry/20180503/1525338512}\\
\url{https://www.overleaf.com/learn/latex/Japanese#The_pTeX_engine}
\end{thebibliography}

\end{document}
